\scanpage{45}
% L1, 3)
3)
\diagram{\begin{tikzcd}[ampersand replacement=\&,column sep=large]
\underline A\arrow[r,shift left,"T"]\&\underline B\arrow[l,shift left,"S"]\&\operatorname{Hom}(x,Sy)\overset{\sim}{\unsure{\mathord{-}|\to}}\operatorname{Hom}(Tx,y)
\end{tikzcd}}
\margin{$\left\{\vphantom{\rule{0pt}{9em}}\right.$}

% L2–3
$x\in\mathrm{Ob}\,\underline A$ conditions équivalentes \uline{si} $S$ \uline{pleinement fidèle}

% L4
\textcircled{a} $x$ est isomorphe à un objet de la forme $Sy$

% L5
\textcircled{b} $x\xrightarrow{\sim}STx$

% L6
\textcircled{c} $\forall\,x'\in A,\ \operatorname{Hom}(x',x)\xrightarrow{\sim}\operatorname{Hom}(Tx',Tx)$

% L7–8
($b\Leftrightarrow c\Rightarrow a$ sans hyp sur $S$, $a\Rightarrow b$ si on suppose $S$ pleinement fidèle).

% L9, 4)
4)
\diagram{\begin{tikzcd}[ampersand replacement=\&,column sep=large]
\underline A\arrow[r,shift left,"T"]\&\underline B\arrow[l,shift left,"S"]\&\operatorname{Hom}(x,Sy)\overset{\sim}{-}\operatorname{Hom}(Tx,y).
\end{tikzcd}}

% L10
Conditions équivalentes

% L11
a) $S$ pleinement fidèle

% L12
b) $TS\overset{\sim}{-}\mathrm{id}_{\underline B}$ isomorphisme

% L13–17
c) $T$ est un foncteur de passage au quotient, i.e. Si on considère \deleted{\unsure{l'ens|l'ens.}} $\unsure{M|\mathcal M}$ des $u\in\mathrm{Fl}\,\underline A$ tels que $T(u)\in\unsure{\mathrm{Flis}|\mathrm{Fliso}|\mathrm{Fl}_{\rm iso}|\mathrm{Isom}}(\underline B)$, et le foncteur $\underline A/M\to\underline B$, ceci est une \uline{équivalence} de catégories.

% archive sheet number: 44
